\section{Testing Simulation and Validation}\label{sec:testing_simulation_and_validation} % (fold)
\subsection{Program Counter (PC)}
\textbf{Testbench code:}
\inputtestbench{test_001/PC_tb.v}

\textbf{Waveform:}
\inputwaveform[0.95]{PC}

\subsection{Address Multiplexer (Address MUX)}
\textbf{Testbench code:}
\inputtestbench{test_002/AM_tb.v}

\textbf{Waveform:}
\inputwaveform[0.95]{AM}

\subsection{Arithmetic Logic Unit (ALU)}
\textbf{Testbench code:}
\inputtestbench{test_003/ALU_tb.v}

\textbf{Waveform:}
\inputwaveform[0.95]{ALU}

\subsection{Controller Unit}
\textbf{Testbench code:}
\inputtestbench{test_004/Controller_tb.v}

\textbf{Waveform:}
\inputwaveform[0.95]{Controller}

\subsection{Instruction Register (IR)}
\textbf{Testbench code:}
\inputtestbench{test_005/IR_tb.v}

\textbf{Waveform:}
\inputwaveform[0.95]{IR}

\subsection{Accumulator Register (AC)}
\textbf{Testbench code:}
\inputtestbench{test_006/AC_tb.v}

\textbf{Waveform:}
\inputwaveform[0.95]{AC}

\subsection{Memory}
\textbf{Testbench code:}
\inputtestbench{test_007/Memory_tb.v}

\textbf{Waveform:}
\inputwaveform[0.95]{Memory}

\subsection{Top-level Integration}
\textbf{Testbench code:}
\inputtestbench{test_008/CPU_tb.v}
\inputtestbench{test_009/CPU2_tb.v}

\textbf{Waveform:}
% \inputwaveform[0.95]{CPU}

% testing_simulation_and_validation (end)

\subsection{Resource, Timing, and Power Evaluation}

In addition to functional verification, the CPU was evaluated at the RTL level in terms of resource usage, timing behavior, and expected power characteristics. The design follows a simple accumulator-based RISC architecture, so the datapath remains compact and includes only the required modules: the Program Counter, Address Multiplexer, Memory, Instruction Register, Accumulator Register, ALU, and Controller.

Resource usage is mainly determined by the storage elements and combinational logic in the RTL design. The processor uses a 5-bit Program Counter, an 8-bit Instruction Register, an 8-bit Accumulator Register, a 3-bit controller state register, and a \(32 \times 8\)-bit memory. The main combinational blocks are the Address MUX, the ALU, and the controller output logic.

Timing behavior is defined by the controller FSM. The controller executes instructions through eight states: \texttt{INST\_ADDR}, \texttt{INST\_FETCH}, \texttt{INST\_LOAD}, \texttt{IDLE}, \texttt{OP\_ADDR}, \texttt{OP\_FETCH}, \texttt{ALU\_OP}, and \texttt{STORE}. Therefore, a normal instruction requires eight clock cycles to complete. With the 10 ns clock period used in simulation, the execution time of a normal instruction is 80 ns. Control-flow instructions such as \texttt{SKZ} and \texttt{JMP} may update the Program Counter, while \texttt{HLT} stops the processor when the halt condition is reached.

\begin{table}[htbp]
\centering
\caption{Resource, Timing, and Power Evaluation}
\label{tab:resource_timing_power}
\small
\renewcommand{\arraystretch}{1.15}
\begin{tabularx}{\textwidth}{@{}p{0.25\textwidth} p{0.23\textwidth} X@{}}
\toprule
\textbf{Metric} & \textbf{Evaluation} & \textbf{Description} \\
\midrule
Instruction width & 8 bits & Includes a 3-bit opcode and a 5-bit operand. \\

Memory organization & \(32 \times 8\) bits & Provides 32 memory locations, each storing 8 bits of data. \\

Datapath width & 8 bits & Used by the ALU, Accumulator, and shared data bus. \\

Address width & 5 bits & Supports addressing of 32 memory locations. \\

Controller states & 8 states & Forms one complete instruction cycle under normal execution. \\

Simulation clock period & 10 ns & Clock period used in the Verilog testbenches. \\

Normal instruction time & 80 ns & Calculated from eight clock cycles with a 10 ns clock period. \\

Sequential resources & PC, IR, AC, FSM, Memory & Store the program address, instruction, accumulator value, control state, and memory contents. \\

Combinational resources & ALU, Address MUX, control logic & Perform arithmetic and logic operations, address selection, and control signal generation. \\

Power evaluation & RTL-level qualitative assessment & Expected to be low due to the small datapath, limited memory size, and simple control logic. \\
\bottomrule
\end{tabularx}
\end{table}


From the evaluation, the CPU satisfies the expected timing behavior of the specified multi-cycle controller. The design completes each normal instruction in eight clock cycles and correctly updates the Program Counter for control-flow instructions. Because the datapath is only 8 bits wide and the memory contains only 32 locations, the hardware complexity remains low. Exact post-synthesis resource utilization and power consumption depend on the target device and synthesis settings; therefore, this report provides an RTL-level evaluation based on the implemented architecture and simulation results.
